\uuid{R6Q0}
\titre{Rang d'une famille de vecteurs/Polynômes}
\auteur{hanani ; mzari}
\organisation{exo7}
\licence{CC BY-NC-SA 4.0}
\datecreate{2020-01-20}
\niveau{L1}
\module{Algèbre}
\chapitre{Espace vectoriel}
\sousChapitre{Base}
\difficulte{2}
\theme{Base et dimension, Rang d'une famille de vecteurs, Polynômes}
\exo7qcmid{lille-2.espaces-vectoriels.033}

\contenu{
\begin{question}
\qclassification{106.04}
\qtags{motcle=Rang d'une famille de vecteurs/Polynômes}
Dans  $\Rr_2[X]$, l'ensemble des polynômes à coefficients réels de degré $\le 2$, on considère les polynômes $P_1= 1-X, P_2= 1+X, P_3 = X^2$ et $P_4=1+X^2$. Quelles sont les assertions vraies ?
\begin{answers}
\bad{Le rang de la famille $\{P_4\}$ est $4$}.
\good{Le rang de la famille $\{P_3,P_4\}$ est $2$}.
\bad{Le rang de la famille $\{P_2,P_3,P_4\}$ est $2$}.
\good{Le rang de la famille $\{P_1,P_2,P_3,P_4\}$ est $3$}.
\end{answers}
\begin{explanations} Le rang d'une famille de vecteurs est la dimension du sous-espace vectoriel
engendré par ces vecteurs. Autrement dit, c'est le nombre maximum de vecteurs de cette famille qui sont linéairement indépendants.
\end{explanations}
\end{question}
}
